%% Packages -----------------------------------------------------------------
\usepackage{mathtools}
\usepackage{booktabs}
\usepackage{longtable}
\usepackage{array}
\usepackage{enumitem}
\usepackage{xurl}
\usepackage{tikz}
\usetikzlibrary{automata,positioning,arrows.meta,decorations.pathreplacing,calc,fit,backgrounds,shapes.geometric,patterns}


%% Lean identifiers: typewriter, with URL-style line breaking.
%% Breaks are allowed only after a dot or an underscore (xurl would break anywhere).
\DeclareUrlCommand\lean{\urlstyle{tt}\def\UrlBreaks{\do\.\do\_}\def\UrlBigBreaks{}\def\UrlNoBreaks{}}

%% Long formulas: allow slightly looser lines instead of overfull boxes.
\setlength{\emergencystretch}{2em}

%% Signature, trees and the order --------------------------------------------
\newcommand{\Bot}{\bot}
\newcommand{\Top}{\top}
\newcommand{\BF}{\{\Bot,f\}}
\newcommand{\BFT}{\{\Bot,f,\Top\}}
\newcommand{\Dir}{[n]}                          % argument positions 1..n
\newcommand{\pol}{\operatorname{pol}}           % polarity of a path
\newcommand{\dom}{\operatorname{dom}}
\newcommand{\Trees}{\mathcal{T}}                % all trees
\newcommand{\TreesReg}{\mathcal{T}_{\mathrm{reg}}}
\newcommand{\TreesFin}{\mathcal{T}_{\mathrm{fin}}}
\newcommand{\Cls}{\mathcal{C}}                  % a class of trees
\newcommand{\fall}{f^{\omega}}                  % the tree with f at every path
\newcommand{\eps}{\varepsilon}
\newcommand{\entails}{\models}

%% Words ----------------------------------------------------------------------
\newcommand{\len}[1]{\lvert #1\rvert}
\newcommand{\wsub}[3]{#1[#2{:}#3]}             % factor w[i:j]
\newcommand{\wpre}[2]{#1[{:}#2]}               % prefix w[:i]
\newcommand{\wsuf}[2]{#1[#2{:}]}               % suffix w[i:]
\newcommand{\pre}{\preceq}                     % prefix order on words
\newcommand{\Alp}{A}                           % alphabet of a reader
\newcommand{\lcsuf}{\operatorname{lcs}}        % length of the longest common suffix (in blocks)

%% Cones, atoms and derivations -------------------------------------------
\newcommand{\Cone}[2]{C^{#1}(#2)}              % cone of sign #1 of the tree #2
\newcommand{\sgp}{{+}}
\newcommand{\sgm}{{-}}
\newcommand{\flip}[1]{\overline{#1}}           % opposite sign
\newcommand{\SV}{V^{\pm}}                      % signed variables
\newcommand{\Atoms}{\mathcal{A}}
\newcommand{\der}{\vdash}                      % derivation by rules
\newcommand{\mv}{\Rightarrow}                  % derivation by moves
\newcommand{\LC}{\mathrm{LC}}                  % least closure
\newcommand{\Hull}{\mathrm{Hull}}
\newcommand{\NegV}{\mathsf{Neg}}               % signed variables with a forbidden root
\newcommand{\FctV}{\mathsf{Fct}}               % signed variables reachable from facts
\newcommand{\Up}{\mathsf{U}}                   % upper run
\newcommand{\Dn}{\mathsf{D}}                   % lower run
\newcommand{\stepU}{\delta^{\uparrow}}
\newcommand{\stepD}{\delta^{\downarrow}}

%% Readers and holes ----------------------------------------------------------
\newcommand{\Rej}{\mathsf{Rej}}
\newcommand{\Adm}{\mathsf{Adm}}
\newcommand{\dead}{\dagger}
\newcommand{\Chan}{\mathsf{Ch}}                % channel
\newcommand{\Vis}{\mathsf{Vis}}                % visited states
\newcommand{\tot}[1]{\lVert #1\rVert}          % total length of a tuple
\newcommand{\dominates}{\sqsubseteq}
\newcommand{\Hole}{\mathsf{H}}                 % small-hole bound
\newcommand{\TailB}{\Lambda}                   % rejected-tail bound
\newcommand{\FreshB}{\mathsf{J}}               % fresh-tuple bound
\newcommand{\expo}{\operatorname{exp}}         % iterated exponential
\newcommand{\cov}{\widehat}                    % cover of a hole

%% Fresh tuples ---------------------------------------------------------------
\newcommand{\mk}[2]{\sigma^{#1}_{#2}}          % marker of level #1 of coordinate #2
\newcommand{\Act}{\mathrm{Act}}
\newcommand{\Anc}{\mathrm{Anc}}
\newcommand{\dcut}{\operatorname{cut}}
\newcommand{\lab}{\operatorname{lab}}
\newcommand{\Marg}{\mathsf{m}}                 % deletion margin

%% Rejected tails ---------------------------------------------------------------
\newcommand{\onset}{\operatorname{on}}
\newcommand{\per}{\operatorname{per}}
\newcommand{\Dict}{\mathcal{D}}
\newcommand{\node}{\operatorname{node}}
\newcommand{\gap}{\operatorname{gap}}
\newcommand{\Lb}{\mathsf{L}}                   % search length for closed factors
\newcommand{\Wb}{\mathsf{B}}                   % window envelope
\newcommand{\Fb}{\Phi}                         % factor length bound
\newcommand{\Qc}{\mathsf{Q}}                   % number of periodic factors: 4+2(k-1)(N-1)
\newcommand{\cmap}{\operatorname{im}}          % cut map (image of a retained cut) of a deletion of positions
\newcommand{\proj}{\operatorname{pr}}          % projection of the cover of a hole

%% Hyperlink targets of numbered statements -----------------------------------
%% The class numbers theorems within sections, but the hyperref anchors use only
%% the number inside the section (definition.1 in Sections 2 and 4, ...), so links
%% to later statements jumped to earlier ones. Include the section in the anchor.
\makeatletter
\@for\pf@env:={theorem,lemma,corollary,proposition,exercise,definition,conjecture,%
  observation,example,note,remark,claim}\do{%
  \@ifundefined{theH\pf@env}{}{%
    \expandafter\def\csname theH\pf@env\endcsname{\theHsection.\arabic{theorem}}}}
\makeatother

%% Lists of references without a serial comma, as in the rest of the text.
\renewcommand{\creflastconjunction}{ and\nobreakspace}
